\documentclass[11pt]{article}

% ---------- Page and fonts ----------
\usepackage[a4paper,margin=1in]{geometry}
\usepackage[utf8]{inputenc}
\usepackage[T1]{fontenc}
\usepackage{microtype}

% ---------- Mathematics ----------
\usepackage{amsmath,amssymb,amsthm,mathtools}
\numberwithin{equation}{section}
\allowdisplaybreaks

% ---------- Lists, links, references ----------
\usepackage{enumitem}
\usepackage[hidelinks]{hyperref}
\usepackage{cleveref}

% ---------- Theorem environments (shared counter, cleveref-aware) ----------
\usepackage{aliascnt}
\newcommand{\sibthm}[3]{%
  \theoremstyle{#3}%
  \newaliascnt{#1}{theorem}%
  \newtheorem{#1}[#1]{#2}%
  \aliascntresetthe{#1}%
  \expandafter\newcommand\csname #1autorefname\endcsname{#2}%
}

\theoremstyle{plain}
\newtheorem{theorem}{Theorem}[section]
\sibthm{lemma}{Lemma}{plain}
\sibthm{proposition}{Proposition}{plain}
\sibthm{corollary}{Corollary}{plain}
\sibthm{conjecture}{Conjecture}{plain}
\sibthm{definition}{Definition}{definition}
\sibthm{example}{Example}{definition}
\sibthm{remark}{Remark}{remark}

\crefname{theorem}{Theorem}{Theorems}
\crefname{lemma}{Lemma}{Lemmas}
\crefname{proposition}{Proposition}{Propositions}
\crefname{corollary}{Corollary}{Corollaries}
\crefname{definition}{Definition}{Definitions}
\crefname{remark}{Remark}{Remarks}
\Crefname{theorem}{Theorem}{Theorems}
\Crefname{lemma}{Lemma}{Lemmas}
\Crefname{proposition}{Proposition}{Propositions}
\Crefname{corollary}{Corollary}{Corollaries}
\Crefname{definition}{Definition}{Definitions}
\Crefname{remark}{Remark}{Remarks}

% ---------- Numbered proof steps ----------
\newlist{steps}{enumerate}{1}
\setlist[steps]{label=\textup{\textbf{Step \arabic*.}}, leftmargin=*,
                labelsep=0.6em, itemsep=0.6ex, topsep=0.6ex}

% ---------- Notation ----------
\newcommand{\dip}{d^{+}}
\newcommand{\Np}{N^{+}}
\newcommand{\chl}{\chi_{\ell}}
\newcommand{\chpl}{\chi_{\ell}'}
\newcommand{\LG}{L}

\title{\bfseries Kernel-Perfect Orientations of Line Graphs\\
of Complete Graphs}
\author{Siddharth Narayan Mishra\\[2pt] Roll Number: 123CS0189}
\date{\today}

\begin{document}
\maketitle

\begin{abstract}
The list-edge-colouring conjecture asserts that $\chi(H)=\chl(H)$ for every line graph $H$.
Galvin proved it for line graphs of bipartite graphs, using the kernel lemma of Bondy,
Boppana and Siegel. We prove that for every even $m\ge 4$ the line graph $\LG(K_m)$ admits
no kernel-perfect orientation whose out-degrees are all at most $\chi'(K_m)-1=m-2$.
Consequently the kernel lemma cannot establish $\chpl(K_m)=\chi'(K_m)$ for any even $m$,
which is precisely the case of the conjecture that remains open for complete graphs.
The proof rests on an observation we call the star lemma: in a kernel-perfect orientation of
a line graph, the edges of $G$ incident with a fixed vertex are linearly ordered. This
converts the question into a numerical one about those orders, and a counting argument then
closes it. The result strengthens a theorem of McDonald, who ruled out the same conclusion
for the restricted family of Galvin orientations. All assertions are proved from first
principles, and every numerical claim has been verified by exhaustive computation in the
stated range.
\end{abstract}

\tableofcontents

\section{Introduction}

All graphs are finite and simple unless stated otherwise. For a graph $G$ we write $V(G)$ and
$E(G)$ for its vertex and edge sets, $d_G(v)$ for the degree of $v$, and $\Delta(G)$ for the
maximum degree. The \emph{line graph} $\LG(G)$ has vertex set $E(G)$, two edges of $G$ being
adjacent in $\LG(G)$ exactly when they share an end.

A \emph{list assignment} for a graph $H$ gives each vertex $v$ a set $C(v)$ of colours. An
\emph{$L$-colouring} is a proper colouring $c$ with $c(v)\in C(v)$ for all $v$. The
\emph{list chromatic number} $\chl(H)$ is the least $k$ such that an $L$-colouring exists
whenever $|C(v)|=k$ for every $v$. Since ordinary colouring is the case of equal lists,
$\chi(H)\le\chl(H)$ always, and the gap can be arbitrarily large in general.

\begin{conjecture}[list-edge-colouring conjecture]\label{conj:lecc}
$\chi(H)=\chl(H)$ for every line graph $H$.
\end{conjecture}

Equivalently, $\chpl(G)=\chi'(G)$ for every multigraph $G$, where $\chi'$ is the chromatic
index and $\chpl$ the list chromatic index. The conjecture has been attributed independently
to several authors and has been open since the 1970s. Galvin proved it for bipartite
multigraphs, thereby settling the Dinitz problem. The tool is the following lemma.

\begin{definition}[kernel]\label{def:kernel}
Let $D$ be a digraph. A set $K\subseteq V(D)$ is a \emph{kernel} of $D$ if
\begin{enumerate}[label=\textup{(\roman*)}]
  \item $K$ is independent in the underlying graph of $D$, and
  \item every $u\notin K$ has an arc $u\to w$ with $w\in K$.
\end{enumerate}
$D$ is \emph{kernel-perfect} if every induced subdigraph of $D$ has a kernel.
\end{definition}

\begin{lemma}[Bondy, Boppana and Siegel]\label{lem:bbs}
Let $H$ be a graph and let $D$ be an orientation of $H$ which is kernel-perfect. If
$|C(v)|\ge\dip_D(v)+1$ for every $v$, then $H$ has an $L$-colouring. In particular
$\chl(H)\le\Delta^{+}(D)+1$.
\end{lemma}

The following notion isolates exactly what \cref{lem:bbs} asks of a line graph.

\begin{definition}[$k$-kernel-orientable]\label{def:orientable}
For an integer $k\ge 1$, a graph $H$ is \emph{$k$-kernel-orientable} if it has a
kernel-perfect orientation $D$ with $\dip_D(x)\le k-1$ for every $x\in V(H)$.
\end{definition}

By \cref{lem:bbs}, if $\LG(G)$ is $\chi'(G)$-kernel-orientable then $\chpl(G)=\chi'(G)$, so
\cref{conj:lecc} holds for $G$. Galvin's theorem is exactly the statement that $\LG(G)$ is
$\chi'(G)$-kernel-orientable whenever $G$ is bipartite. Our main result says that this route
is unavailable for complete graphs of even order.

\begin{theorem}\label{thm:main}
Let $m\ge 4$ be even. Then $\LG(K_m)$ is \emph{not} $(m-1)$-kernel-orientable. Since
$\chi'(K_m)=m-1$, the kernel lemma \cref{lem:bbs} cannot establish
$\chpl(K_m)=\chi'(K_m)$ for any even $m$.
\end{theorem}

The restriction to even $m$ is not a weakness of the method but the point of the result.
For odd $m$ one has $\chi'(K_m)=m$, and Haggkvist and Janssen proved $\chpl(K_m)\le m$, so
\cref{conj:lecc} is settled for odd complete graphs. The even case is the one that remains
open, and \cref{thm:main} says the chapter's own method is not available there.

\Cref{thm:main} strengthens a theorem of McDonald, who proved the same conclusion for the
restricted family of \emph{Galvin orientations}. We discuss the relationship in
\cref{sec:mcdonald}.

\section{The star lemma}\label{sec:star}

Throughout this section $G$ is a graph and $D$ is an orientation of $\LG(G)$. For
$v\in V(G)$ let $E(v)$ denote the set of edges of $G$ incident with $v$. Since any two
members of $E(v)$ share the end $v$, the set $E(v)$ is a clique of $\LG(G)$, so $D$ induces a
tournament on it.

\begin{lemma}[star lemma]\label{lem:star}
Let $D$ be a kernel-perfect orientation of $\LG(G)$. Then for every $v\in V(G)$ the tournament
induced by $D$ on $E(v)$ is transitive.
\end{lemma}

\begin{proof}
\begin{steps}
\item Let $T$ be a tournament and let $K$ be a kernel of $T$. Since $K$ is independent and
any two vertices of a tournament are adjacent, $|K|\le 1$; and $K\ne\emptyset$, since
otherwise condition (ii) of \cref{def:kernel} would ask a vertex to have an arc into the
empty set. So $K=\{x\}$ for a single vertex $x$, and (ii) says every other vertex has an arc
to $x$. Hence $x$ has out-degree $0$ in $T$, that is, $x$ is a sink.

\item A directed triangle has no sink, since each of its three vertices has out-degree $1$.
By Step 1 it therefore has no kernel.

\item $E(v)$ is a clique of $\LG(G)$, so every $S\subseteq E(v)$ induces a tournament in $D$.
As $D$ is kernel-perfect, that tournament has a kernel. By Step 2 no such $S$ induces a
directed triangle, so $D$ has no directed triangle inside $E(v)$.

\item A tournament with no directed triangle is transitive. Indeed, suppose $a\to b$ and
$b\to c$ with $a\ne c$. The arc between $a$ and $c$ cannot be $c\to a$, for then
$a\to b\to c\to a$ would be a directed triangle. Hence $a\to c$. \qedhere
\end{steps}
\end{proof}

\begin{remark}\label{rem:bkw-clique}
\Cref{lem:star} is not new in content. Borodin, Kostochka and Woodall proved that an
orientation of a line graph is kernel-perfect if and only if every clique has a kernel and
every directed odd cycle has a chord or a pseudochord. The necessity of the clique condition,
applied to every subset of $E(v)$, yields \cref{lem:star} at once. We have given the direct
argument above so that this article remains self-contained, and because the formulation as a
linear order, rather than as the existence of a kernel, is what makes
\cref{def:rank} available.
\end{remark}

\begin{remark}\label{rem:galvin-shape}
Galvin orients $\LG(K_{n,n})$ by the entries of a Latin square, running along each row from
the smaller entry to the larger and along each column in the opposite direction. Rows and
columns are precisely the sets $E(v)$ for $K_{n,n}$, so his construction is a linear order on
each $E(v)$. \Cref{lem:star} says this shape is forced: no other kind of orientation can be
kernel-perfect.
\end{remark}

\Cref{lem:star} lets us replace an orientation by a family of numbers.

\begin{definition}[the rank function]\label{def:rank}
Let $D$ be a kernel-perfect orientation of $\LG(G)$. For $v\in V(G)$ and $e\in E(v)$ put
\[
  M(v,e)\;=\;\bigl|\{\,f\in E(v)\;:\;e\to f \text{ in } D\,\}\bigr|.
\]
For an edge $e=uv$ of $G$ we abbreviate $M(u,v):=M(u,e)$ and $M(v,u):=M(v,e)$.
\end{definition}

\begin{lemma}\label{lem:rank}
Let $D$ be a kernel-perfect orientation of $\LG(G)$ and let $v\in V(G)$. Then
\begin{enumerate}[label=\textup{(\alph*)}]
  \item $M(v,\cdot)$ is a bijection from $E(v)$ onto $\{0,1,\dots,d_G(v)-1\}$;
  \item for $e,f\in E(v)$ with $e\ne f$ we have $e\to f$ in $D$ if and only if
        $M(v,e)>M(v,f)$;
  \item for every edge $e=uv$ of $G$,
        $\ \dip_D(e)=M(u,v)+M(v,u)$.
\end{enumerate}
\end{lemma}

\begin{proof}
\begin{steps}
\item By \cref{lem:star} the tournament on $E(v)$ is transitive, so it is a linear order:
write $E(v)=\{e_1,\dots,e_d\}$ with $e_i\to e_j$ exactly when $i<j$, where $d=d_G(v)$. Then
$M(v,e_i)=|\{j:j>i\}|=d-i$, which runs over $\{0,1,\dots,d-1\}$ bijectively as $i$ runs from
$1$ to $d$. This proves (a), and (b) follows since $i<j$ is equivalent to $d-i>d-j$.

\item The neighbours of $e=uv$ in $\LG(G)$ are the edges of $G$ sharing $u$ with $e$ and the
edges sharing $v$ with $e$. As $G$ is simple, no edge other than $e$ lies in both $E(u)$ and
$E(v)$, so these two sets of neighbours are disjoint. The out-neighbours of $e$ therefore
split accordingly, giving $\dip_D(e)=M(u,e)+M(v,e)$, which is (c). \qedhere
\end{steps}
\end{proof}

\section{The budget is tight on regular class one graphs}\label{sec:budget}

Recall that $G$ is \emph{class one} if $\chi'(G)=\Delta(G)$.

\begin{lemma}\label{lem:budget}
Let $G$ be $\Delta$-regular and class one, and suppose $\LG(G)$ is
$\chi'(G)$-kernel-orientable, witnessed by $D$. Then
\[
  M(u,v)+M(v,u)\;=\;\Delta-1\qquad\text{for every edge } uv \text{ of } G .
\]
\end{lemma}

\begin{proof}
\begin{steps}
\item Let $n=|V(G)|$, so $|E(G)|=n\Delta/2$. Two edges of $G$ are adjacent in $\LG(G)$
exactly when they share an end, and by simplicity they share at most one end, so each
vertex $v$ contributes exactly $\binom{d_G(v)}{2}$ edges to $\LG(G)$ and no edge of $\LG(G)$
is counted twice. Hence
\[
  |E(\LG(G))|\;=\;\sum_{v\in V(G)}\binom{d_G(v)}{2}\;=\;n\binom{\Delta}{2}
  \;=\;\frac{n\Delta(\Delta-1)}{2}.
\]

\item Every edge of $\LG(G)$ receives exactly one orientation in $D$, so
$\sum_{e\in E(G)}\dip_D(e)=|E(\LG(G))|$.

\item By hypothesis $\dip_D(e)\le\chi'(G)-1=\Delta-1$ for every $e$, and there are
$|E(G)|=n\Delta/2$ edges, so
\[
  \sum_{e\in E(G)}\dip_D(e)\;\le\;(\Delta-1)\cdot\frac{n\Delta}{2}
  \;=\;\frac{n\Delta(\Delta-1)}{2}.
\]

\item Steps 1 to 3 give equality throughout, so $\dip_D(e)=\Delta-1$ for every edge $e$.
\Cref{lem:rank}(c) now gives the claim. \qedhere
\end{steps}
\end{proof}

\section{Proof of the main theorem}\label{sec:main}

\begin{proof}[Proof of \cref{thm:main}]
Let $m\ge 4$ be even, write $G=K_m$, and put $s=m-2$. Then $G$ is $(m-1)$-regular and
$\chi'(G)=m-1$, so $G$ is class one with $\Delta=m-1$ and $s=\Delta-1$.

Suppose for a contradiction that $\LG(K_m)$ is $(m-1)$-kernel-orientable, witnessed by a
kernel-perfect orientation $D$ with all out-degrees at most $m-2$. Let $M$ be the rank
function of \cref{def:rank}.

\begin{steps}
\item By \cref{lem:budget},
\begin{equation}\label{eq:tight}
  M(u,v)+M(v,u)=s\qquad\text{for every edge } uv .
\end{equation}
By \cref{lem:rank}(a), each $M(u,\cdot)$ is a bijection onto $\{0,1,\dots,s\}$, since
$d_G(u)=m-1=s+1$.

\item Fix any vertex $i$ of $K_m$. By \cref{lem:rank}(a) there is a unique neighbour $j$ of
$i$ with $M(i,j)=0$, and then \eqref{eq:tight} gives $M(j,i)=s$.

\item Let $k$ be any vertex other than $i$ and $j$; there are exactly $m-2=s$ such $k$. Since
$K_m$ is complete, $ij$, $jk$ and $ik$ are edges of $K_m$ and hence vertices of $\LG(K_m)$,
pairwise adjacent, so they induce a triangle in $\LG(K_m)$. We determine its orientation.

\item \emph{The arc between $ij$ and $jk$.} These two share the end $j$, so by
\cref{lem:rank}(b), $ij\to jk$ if and only if $M(j,i)>M(j,k)$. Now $M(j,i)=s$ is the largest
value in the range of $M(j,\cdot)$, and $M(j,k)\ne M(j,i)$ because $M(j,\cdot)$ is injective
and $k\ne i$. Hence $M(j,i)>M(j,k)$ and
\[
  ij\longrightarrow jk .
\]

\item \emph{The arc between $ik$ and $ij$.} These share the end $i$, so $ik\to ij$ if and only
if $M(i,k)>M(i,j)=0$. Since $M(i,\cdot)$ is injective and $k\ne j$, we have $M(i,k)\ne 0$,
hence $M(i,k)>0$ and
\[
  ik\longrightarrow ij .
\]

\item \emph{The arc between $jk$ and $ik$.} These share the end $k$, so $jk\to ik$ if and only
if $M(k,j)>M(k,i)$. By \eqref{eq:tight}, $M(k,j)=s-M(j,k)$ and $M(k,i)=s-M(i,k)$, so this
condition is equivalent to $M(i,k)>M(j,k)$.

\item If $M(i,k)>M(j,k)$ then, combining Steps 4, 5 and 6, we obtain
\[
  ij\longrightarrow jk\longrightarrow ik\longrightarrow ij ,
\]
a directed triangle in $D$. By Step 2 of the proof of \cref{lem:star} a directed triangle has
no kernel, and it is an induced subdigraph of $D$, contradicting kernel-perfectness. Hence
$M(i,k)\le M(j,k)$. Moreover $M(k,i)\ne M(k,j)$ because $M(k,\cdot)$ is injective and
$i\ne j$, so by \eqref{eq:tight} $M(i,k)\ne M(j,k)$. Therefore
\begin{equation}\label{eq:key}
  M(i,k)<M(j,k)\qquad\text{for every } k\notin\{i,j\}.
\end{equation}

\item Sum \eqref{eq:key} over the $s$ vertices $k\notin\{i,j\}$. Since $M(i,\cdot)$ is a
bijection onto $\{0,1,\dots,s\}$ and omits only the value $M(i,j)=0$ when $k$ ranges over
$V\setminus\{i,j\}$,
\[
  \sum_{k\notin\{i,j\}}M(i,k)\;=\;\Bigl(\sum_{t=0}^{s}t\Bigr)-M(i,j)
  \;=\;\frac{s(s+1)}{2}-0\;=\;\frac{s(s+1)}{2}.
\]
Likewise $M(j,\cdot)$ omits only the value $M(j,i)=s$, so
\[
  \sum_{k\notin\{i,j\}}M(j,k)\;=\;\frac{s(s+1)}{2}-s .
\]

\item Summing \eqref{eq:key} over all $s\ge 1$ choices of $k$ gives a strict inequality:
\[
  \frac{s(s+1)}{2}\;<\;\frac{s(s+1)}{2}-s ,
\]
that is $0<-s$. Since $s=m-2\ge 2$ this is false. The assumed orientation cannot exist. \qedhere
\end{steps}
\end{proof}

\begin{remark}\label{rem:where-force}
The argument never analyses a triangle in isolation. \Cref{prop:cubic} below shows that for
$\Delta\ge 4$ a single triangle imposes no contradiction at all, so no local argument can
prove \cref{thm:main}. The force comes from choosing $j$ at the extreme of the order at $i$,
which fixes two of the three arcs of \emph{every} triangle through $ij$ at once and makes all
of them yield the same inequality \eqref{eq:key}. The contradiction is then a statement about
row sums.
\end{remark}

\section{Sharpness and scope}\label{sec:scope}

Two remarks delimit \cref{thm:main}. The first shows that at $\Delta=3$ a purely local
argument does suffice, and the second shows that it does not for larger $\Delta$.

\begin{proposition}\label{prop:cubic}
Let $G$ be $3$-regular and class one. If $G$ contains a triangle, then $\LG(G)$ is not
$3$-kernel-orientable.
\end{proposition}

\begin{proof}
Suppose $D$ witnesses $3$-kernel-orientability. By \cref{lem:budget} with $\Delta=3$ we have
$s=2$ and $M(u,v)+M(v,u)=2$ for every edge. Let $ijk$ be a triangle of $G$ and set
$p=M(i,j)$, $q=M(j,k)$ and $r=M(i,k)$, so that $M(j,i)=2-p$, $M(k,j)=2-q$ and $M(k,i)=2-r$.
Injectivity of $M(i,\cdot)$, $M(j,\cdot)$ and $M(k,\cdot)$ gives respectively
\[
  p\ne r,\qquad p+q\ne 2,\qquad q\ne r .
\]
By \cref{lem:rank}(b), as in Steps 4 to 6 above, the triangle $ij,jk,ik$ of $\LG(G)$ is
directed precisely when
\[
  \bigl(p+q<2 \text{ and } r>p \text{ and } r>q\bigr)\quad\text{or}\quad
  \bigl(p+q>2 \text{ and } r<p \text{ and } r<q\bigr).
\]
There are nine pairs $(p,q)\in\{0,1,2\}^2$. The three with $p+q=2$ are excluded. For the
remaining six, avoiding a directed triangle forces $r$ into a value already excluded:
\begin{center}
\begin{tabular}{ccl}
\hline
$(p,q)$ & requirement for no directed triangle & conflict\\
\hline
$(0,0)$ & $r\le 0$ & contradicts $r\ne p=0$\\
$(0,1)$ & $r\le 1$, and $r\ne 0$, so $r=1$ & contradicts $r\ne q=1$\\
$(1,0)$ & $r\ne 2$, and $r\ne 1$, so $r=0$ & contradicts $r\ne q=0$\\
$(1,2)$ & $r\ne 0$, and $r\ne 1$, so $r=2$ & contradicts $r\ne q=2$\\
$(2,1)$ & $r\ne 0$, and $r\ne 2$, so $r=1$ & contradicts $r\ne q=1$\\
$(2,2)$ & $r=2$ & contradicts $r\ne p=2$\\
\hline
\end{tabular}
\end{center}
Every case fails, so the triangle of $\LG(G)$ is directed and has no kernel.
\end{proof}

\begin{remark}\label{rem:not-triangles}
\Cref{prop:cubic} does not extend to larger degrees, in two distinct senses. First, the local
obstruction disappears: for $s=\Delta-1\ge 3$ there are triples $(p,q,r)$ satisfying the three
injectivity conditions and yielding an undirected triangle, for instance $(p,q,r)=(0,2,1)$.
Exhaustive enumeration gives $0,0,6,18,42,78,132$ admissible triples for
$s=1,2,3,4,5,6,7$ respectively. Second, the presence of a triangle is not by itself an
obstruction: the circulant $C_8(1,2)$ is $4$-regular and class one and contains eight
triangles, yet admits a system of orders satisfying \eqref{eq:tight} with no directed
triangle.
\end{remark}

\begin{remark}\label{rem:reach}
The proof of \cref{thm:main} uses only that $i$ and $j$ are adjacent and that every other
vertex is adjacent to both. For a $\Delta$-regular class one graph this says
$N(i)\setminus\{j\}\subseteq N(j)$, and since both sets have size $\Delta-1$ it forces
$N[i]=N[j]$. Because $j$ is determined by $D$ rather than chosen, the argument as stated
applies exactly to disjoint unions of complete graphs of even order. Extending it further
requires a different way of aggregating the inequalities \eqref{eq:key}.
\end{remark}

\section{Computational verification}\label{sec:computation}

Every numerical assertion above has been checked independently.

\begin{enumerate}[label=\textup{(\arabic*)}]
  \item \emph{$K_4$, by exhaustive orientation search.} $\LG(K_4)$ is the octahedron, with
  $6$ vertices and $12$ edges. All $2^{12}=4096$ orientations were enumerated; the $38$ with
  all out-degrees at most $2$ were tested for kernel-perfectness by examining every one of the
  $2^{6}-1$ induced subdigraphs. Every one of the $38$ contains a directed triangle, so none
  is kernel-perfect. This confirms \cref{thm:main} for $m=4$ without using
  \cref{lem:star}.

  \item \emph{$K_4$, $K_6$ and $K_8$, by exhaustive search over rank functions.} Backtracking
  over all $M$ satisfying \eqref{eq:tight}, injectivity of every row, and absence of a
  directed triangle finds no solution in each case, agreeing with \cref{thm:main}.

  \item \emph{\Cref{lem:star}.} For every tournament on at most $6$ vertices, the properties
  ``every non-empty subset contains a sink'' and ``transitive'' select the same tournaments,
  namely $2$, $8$, $64$, $1024$ and $32768$ of them on $2,3,4,5,6$ vertices.

  \item \emph{\Cref{rem:not-triangles}.} The counts $0,0,6,18,42,78,132$ and the example
  $C_8(1,2)$ were produced by direct enumeration.
\end{enumerate}

\section{Relation to McDonald's theorem}\label{sec:mcdonald}

McDonald defines a \emph{Galvin orientation} of $\LG(G)$ from two pieces of data: a partition
of $V(G)$ into sets $U$ and $D$, and a single $k$-edge-colouring $\varphi$ of $G$. For
adjacent $e_1,e_2\in E(G)$ with $\varphi(e_1)<\varphi(e_2)$ sharing a vertex $w$, the edge of
$\LG(G)$ between them is oriented from $e_2$ to $e_1$ if $w\in D$, and from $e_1$ to $e_2$ if
$w\in U$. Such an orientation is \emph{proper with respect to $k$} if it is kernel-perfect and
has all out-degrees at most $k-1$. She proves:

\begin{theorem}[McDonald]\label{thm:mcdonald}
$\LG(K_n)$ has no proper Galvin orientation with respect to $\chi'(K_n)$, for any $n\ge 4$.
\end{theorem}

\Cref{thm:main} is a strengthening of \cref{thm:mcdonald} in the even case, for the following
reason. In a Galvin orientation every vertex orders its incident edges by the \emph{same}
colouring $\varphi$, choosing only whether to read it upwards or downwards. By
\cref{lem:star}, a general kernel-perfect orientation gives each vertex an arbitrary linear
order on its incident edges. To recover a single global colouring from a rank function $M$
one would need to set $\varphi(uv)=M(v,u)+1$ consistently, which requires each vertex to play
the same role for all of its edges, that is, requires a partition of $V(G)$ with one end of
every edge in each part. That is precisely bipartiteness. For non-bipartite $G$, and in
particular for $K_m$, the rank functions covered by \cref{thm:main} therefore form a strictly
larger family than the Galvin orientations covered by \cref{thm:mcdonald}.

Informally, \cref{thm:mcdonald} says that Galvin's construction does not reach complete
graphs, while \cref{thm:main} says that the kernel lemma does not, however the orientation is
built.

McDonald also proves a positive companion result: $\LG(K_{\Delta+1})$ does admit a proper
Galvin orientation with respect to roughly $\lfloor 3\Delta/2\rfloor$, which is far from the
bound $\chpl(K_n)\le n$ of Haggkvist and Janssen.

\begin{remark}\label{rem:status}
The two places where \cref{thm:main} would most plausibly already appear have been consulted.
The article \emph{On kernel-perfect orientations of line graphs} by Borodin, Kostochka and
Woodall, Discrete Math.\ \textbf{191} (1998) 45--49, proves only the characterisation quoted in
\cref{rem:bkw-clique}; its statement concerns cliques and odd cycles in general and makes no
assertion about complete graphs or about out-degree bounds. McDonald's article proves
\cref{thm:mcdonald}, which as explained above concerns the strictly smaller family of Galvin
orientations. Neither contains \cref{thm:main}.

This is evidence and not proof. The author has not read the full text of the 1998 article,
only its abstract and the statement as quoted by McDonald, and a search of the literature
cannot establish a negative. \Cref{thm:main} should be regarded as a strengthening of
\cref{thm:mcdonald} whose novelty is likely but unconfirmed.
\end{remark}

\section{Concluding remarks}

\Cref{thm:main} closes the kernel route to \cref{conj:lecc} for complete graphs of even order,
which is exactly where the conjecture remains open for this family. It is consistent with the
shape of the literature: the positive results on complete graphs, such as those of Haggkvist
and Janssen, of Schauz for $K_{p+1}$ with $p$ an odd prime, and the recent families
$K_{p-1}$ and $K_{2p}$, all proceed through the Alon and Tarsi polynomial method and its
formulation in terms of signed one-factorizations, rather than through kernels.

It is worth recording that the two methods are incomparable rather than nested. The Alon and
Tarsi method proves $\chpl(K_4)=3$, which by \cref{thm:main} the kernel lemma cannot. In the
other direction, Galvin's orientation of $\LG(K_{3,3})$ is kernel-perfect with all out-degrees
equal to $2$, so the kernel lemma proves $\chpl(K_{3,3})=3$, whereas the Alon and Tarsi
certificate for that orientation vanishes and the Alon and Tarsi number of $\LG(K_{3,3})$ is
$4$. Neither method subsumes the other.

\end{document}
